% ===========================================================================
%  comune/stile-codice.tex — come si mostra il codice
% ===========================================================================
%
%  PERCHE' listings E NON minted. minted colora meglio, ma per farlo chiama
%  Pygments durante la compilazione: serve Python installato e il permesso
%  --shell-escape, che molti editor LaTeX negano per difetto e che le
%  compilazioni automatiche rifiutano. Una documentazione che non compila sulla
%  macchina di chi la legge non e' documentazione. listings sta dentro LaTeX e
%  compila ovunque, e la differenza di resa non vale quel prezzo.
%
%  Tre modi di mostrare codice, e vanno usati per cose diverse:
%
%    \begin{codice}[python]{didascalia}{percorso/file.py}  blocco con cornice
%    \inputcodice{python}{didascalia}{file.tex}{percorso}  da file esterno
%    \cod{nome_simbolo}                                     dentro una frase
%
%  Il PERCORSO non e' decorazione: chi legge la documentazione ha il sorgente
%  aperto accanto, e un frammento senza il suo file costringe a cercarlo.
% ===========================================================================
\RequirePackage{listings}
\RequirePackage{tcolorbox}
\tcbuselibrary{listings, skins, breakable}

% --- stile comune ----------------------------------------------------------
\lstdefinestyle{lyraBase}{
  basicstyle      = \ttfamily\footnotesize\color{lyraInchiostro},
  keywordstyle    = \color{codParola}\bfseries,
  commentstyle    = \color{codCommento}\itshape,
  stringstyle     = \color{codStringa},
  numberstyle     = \ttfamily\scriptsize\color{codRiga},
  numbers         = left,
  numbersep       = 9pt,
  stepnumber      = 1,
  showstringspaces= false,
  breaklines      = true,
  breakatwhitespace = false,
  % La freccia dice che la riga continua: senza, una riga spezzata si legge
  % come due istruzioni, ed e' il tipo di errore che si propaga nel copia-incolla.
  prebreak        = \raisebox{0ex}[0ex][0ex]{\color{lyraAccento}\ensuremath{\hookleftarrow}},
  tabsize         = 4,
  columns         = fullflexible,
  keepspaces      = true,
  upquote         = true,
  extendedchars   = true,
  literate        =
    {à}{{\`a}}1 {è}{{\`e}}1 {é}{{\'e}}1 {ì}{{\`\i}}1 {ò}{{\`o}}1 {ù}{{\`u}}1
    {À}{{\`A}}1 {È}{{\`E}}1 {Ì}{{\`I}}1 {Ò}{{\`O}}1 {Ù}{{\`U}}1
    {€}{{\texteuro}}1 {°}{{\textdegree}}1 {·}{{\textperiodcentered}}1
    {—}{{---}}1 {–}{{--}}1 {«}{{\guillemotleft}}1 {»}{{\guillemotright}}1
    {“}{{\textquotedblleft}}1 {”}{{\textquotedblright}}1
    {’}{{\textquoteright}}1 {‘}{{\textquoteleft}}1
    {→}{{$\rightarrow$}}1 {←}{{$\leftarrow$}}1 {✓}{{$\checkmark$}}1,
}

% ---------------------------------------------------------------------------
%  REGOLA: dentro un blocco di codice si scrive SOLO ASCII.
%
%  Niente lettere accentate, niente caratteri di disegno, niente frecce.
%  `literate` non basta, e nemmeno il pacchetto listingsutf8: tcblisting
%  scrive il contenuto in un file intermedio e lo rilegge, quindi i byte
%  arrivano al kernel prima che listings possa intercettarli. L'errore e'
%  «Invalid UTF-8 byte sequence», segnalato nel file .listing e senza dire
%  quale blocco l'ha causato.
%
%  Quando un valore del sorgente contiene un accento, si translittera e si
%  mette un commento accanto. Gli alberi di cartelle si disegnano in ASCII
%  (|-- e `--), oppure in TikZ, dove non passano da listings.
%  `make controlla` segnala le violazioni prima che fermino la compilazione.
% ---------------------------------------------------------------------------

% --- i linguaggi che compaiono nel progetto ---------------------------------
\lstdefinelanguage{JavaScriptLyra}{
  keywords = {break, case, catch, class, const, continue, debugger, default,
    delete, do, else, export, extends, finally, for, function, if, import, in,
    instanceof, let, new, of, return, super, switch, this, throw, try, typeof,
    var, void, while, with, yield, async, await, static, get, set},
  keywordstyle = \color{codParola}\bfseries,
  ndkeywords = {true, false, null, undefined, document, window, console,
    fetch, JSON, Promise, Array, Object, Math},
  ndkeywordstyle = \color{lyraInchiostro!75}\bfseries,
  sensitive = true,
  comment = [l]{//},
  morecomment = [s]{/*}{*/},
  morestring = [b]',
  morestring = [b]",
  morestring = [b]`,
}

\lstdefinelanguage{TypeScriptLyra}{
  morekeywords = [1]{break, case, catch, class, const, continue, default,
    delete, do, else, enum, export, extends, finally, for, function, if,
    implements, import, in, instanceof, interface, let, new, of, private,
    protected, public, readonly, return, super, switch, this, throw, try,
    type, typeof, var, void, while, yield, async, await, as, satisfies},
  morekeywords = [2]{string, number, boolean, any, unknown, never, Promise,
    Record, Array, Partial, Required},
  keywordstyle = [1]\color{codParola}\bfseries,
  keywordstyle = [2]\color{lyraInchiostro!70},
  sensitive = true,
  comment = [l]{//},
  morecomment = [s]{/*}{*/},
  morestring = [b]',
  morestring = [b]",
  morestring = [b]`,
}

\lstdefinelanguage{Console}{
  basicstyle = \ttfamily\footnotesize,
  morecomment = [l]{\#},
  sensitive = false,
}

% --- l'ambiente principale --------------------------------------------------
%  #1 linguaggio, con il nome CORTO (python, js, ts, sql, json, yaml, html,
%     css, console): \lstlingua lo traduce nel nome che listings conosce, così
%     chi scrive non deve ricordare che JSON non esiste e che gli somiglia Java
%  #2 didascalia
%  #3 percorso del file nel progetto (vuoto = nessuno)
\newtcblisting{codiceInterno}[3]{
  listing only,
  % Sia lo stile sia la lingua devono stare QUI dentro: `listing options`
  % sostituisce `listing style` invece di aggiungersi, e il sintomo e' che
  % meta' delle impostazioni (breaklines, numeri di riga) sparisce senza
  % dirlo. Le righe lunghe escono dal riquadro, ed e' l'unico segno.
  listing options = {style = lyraBase, language = #1},
  breakable,
  enhanced,
  colback   = codFondo,
  colframe  = codBordo,
  boxrule   = 0.5pt,
  arc       = 3pt,
  left      = 3pt, right = 3pt, top = 4pt, bottom = 4pt,
  toptitle  = 2pt, bottomtitle = 2pt,
  title     = {\footnotesize\sffamily\bfseries\color{lyraInchiostro}#2%
               \ifstrempty{#3}{}{\hfill\normalfont\ttfamily\scriptsize
                 \color{lyraInchiostro!60}#3}},
  colbacktitle = lyraPannello!60!white,
  coltitle  = lyraInchiostro,
  fonttitle = \footnotesize,
  titlerule = 0pt,
  before skip = 10pt, after skip = 10pt,
}

% Nomi corti per i linguaggi, cosi' chi scrive non deve ricordare i nomi
% interni di listings ne' che JSON e YAML non esistono.
%
% Si usa \lstalias e non una macro che traduce: l'opzione `language` di
% listings viene letta senza espandere, quindi una macro li' dentro produce
% «Missing \endcsname» alla fine del blocco di codice — errore che non dice
% niente di cio' che e' successo.
\lstalias{console}{Console}
\lstalias{python}{Python}
\lstalias{js}{JavaScriptLyra}
\lstalias{ts}{TypeScriptLyra}
\lstalias{sql}{SQL}
\lstalias{json}{Java}      % JSON non esiste in listings: Java gli somiglia
\lstalias{yaml}{Python}    % idem per YAML: commenti con #, due punti, indentazione
\lstalias{html}{HTML}
\lstalias{css}{HTML}
\lstalias{testo}{}

% --- codice da file esterno -------------------------------------------------
%  #1 linguaggio  #2 didascalia  #3 percorso del file su disco  #4 percorso da
%  mostrare. Si usa per gli estratti lunghi tenuti fuori dal testo.
\newcommand{\inputcodice}[4]{%
  \begin{tcolorbox}[
    breakable, enhanced,
    colback = codFondo, colframe = codBordo, boxrule = 0.5pt, arc = 3pt,
    left = 3pt, right = 3pt, top = 4pt, bottom = 4pt,
    title = {\footnotesize\sffamily\bfseries\color{lyraInchiostro}#2%
             \hfill\normalfont\ttfamily\scriptsize\color{lyraInchiostro!60}#4},
    colbacktitle = lyraPannello!60!white, coltitle = lyraInchiostro,
    fonttitle = \footnotesize, titlerule = 0pt]
    \lstinputlisting[style=lyraBase, language=#1]{#3}
  \end{tcolorbox}}

% --- codice dentro una frase ------------------------------------------------
%  Fondo tenue invece del grassetto: dentro un paragrafo il grassetto chiama
%  l'occhio e spezza la lettura, mentre un nome di funzione va riconosciuto
%  senza essere gridato.
%  Le due penalita' attorno alla scatola non sono decorazione. Un \tcbox e'
%  una scatola indivisibile: se capita dopo una parentesi o un trattino --
%  cioe' senza spazio davanti -- TeX non ha dove spezzare la riga e la lascia
%  sbordare nel margine. Sono state la causa di quasi tutte le righe troppo
%  lunghe della prima stesura. Con queste, la scatola puo' sempre andare a
%  capo per conto suo.
%
%  Resta un limite: una scatola piu' larga della giustezza sborda comunque.
%  Per un percorso lungo si usa \file{} (testo normale, si spezza) oppure lo
%  si mette in un blocco di codice.
\DeclareRobustCommand{\cod}[1]{%
  \penalty50
  \tcbox[on line, boxsep = 0.6pt, left = 1.4pt, right = 1.4pt, top = 0pt,
         bottom = 0pt, arc = 2pt, boxrule = 0pt,
         colback = lyraPannello, colframe = lyraPannello]
        {\ttfamily\small #1}%
  \penalty0 }

% Percorso di file nel testo.
%  Usa lo stesso ciclo di \percorsoHTTP (definito in comune/macro.tex): un nome
%  come `diagnose_regime_signals.py' e' lungo quattro centimetri e senza punti
%  di sillabazione, quindi in una colonna stretta sborda. Con i punti di
%  interruzione dopo / _ . e - va a capo dove si legge comunque.
%  Il ciclo richiede ASCII e niente spazi: per il resto c'e' \cod{}.
\DeclareRobustCommand{\file}[1]{{\color{lyraInchiostro!78}\percorsoHTTP{#1}}}

% ===========================================================================
%  Un COMMENTO del sorgente citato per esteso
% ===========================================================================
%  Diversa da `codiceInterno`, e la differenza non e' estetica. Molti commenti
%  di questo progetto sono prose italiane di dieci righe che spiegano una
%  decisione: sono la parte piu' utile da citare, e dentro un blocco di
%  listings non ci stanno — listings vuole ASCII puro (vedi sopra), e
%  translitterare «perche'» e le lettere accentate significa citare qualcosa
%  che nel sorgente non c'e'.
%
%  Qui il testo e' testo LaTeX normale: accenti, virgolette e trattini lunghi
%  si compongono come nel resto del documento, e la citazione resta fedele.
%
%    \begin{commento}{percorso/del/file.py:120}
%    Il testo del commento, a capo come nel sorgente.
%    \end{commento}
%
%  Il corsivo e il filetto a sinistra dicono «questo l'ha scritto chi ha
%  scritto il codice, non chi ha scritto la guida».
\newtcolorbox{commento}[1]{
  enhanced, breakable,
  colback   = codFondo, colframe = lyraInchiostro!30,
  boxrule   = 0pt, leftrule = 2.2pt, arc = 2pt,
  left = 9pt, right = 9pt, top = 6pt, bottom = 6pt,
  fontupper = \small\itshape,
  title     = {\normalfont\ttfamily\scriptsize\color{lyraInchiostro!60}#1},
  colbacktitle = codFondo, coltitle = lyraInchiostro,
  titlerule = 0pt, toptitle = 2pt, bottomtitle = 0pt,
  before skip = 10pt, after skip = 10pt,
}
